\bookchapter{第五十八章：クイックリファレンスとリソース}

\chapterintro{この章は、単独で使えるクイックリファレンスです。書き方を思い出せないときや、素早く決めたいときは、直接ここを開いてください。「ポケットカード」として使うのがおすすめです。よく使うコードの速査、よくある質問、学習リソースを、1ページも漏らさずまとめました。}

\booksection{よく使うコードのクイックカード}

\begin{center}
\begin{tabularx}{\linewidth}{|X|X|}
\hline
\textbf{やりたいこと} & \textbf{キーコード} \tabularnewline
\hline
可変状態を1つ覚えておく & var x by remember \{ mutableStateOf(初期値) \} \tabularnewline
画面回転でも状態を失わない & rememberSaveable \{ mutableStateOf(初期値) \} \tabularnewline
縦並び／横並びのレイアウト & \texttt{Column \{ \}}/\texttt{Row \{ \}} \tabularnewline
遅延読み込みリスト & LazyColumn \{ items(リスト) \{ ... \} \} \tabularnewline
画面遷移 & navController.navigate("ルート") \tabularnewline
コルーチンで待ち時間の仕事をする & \texttt{viewModelScope.\allowbreak launch \{ .\allowbreak .\allowbreak .\allowbreak  \}} \tabularnewline
データフローを監視する & \texttt{val list by flow.\allowbreak collectAsState()} \tabularnewline
小さな設定を保存する & context.dataStore.edit \{ it[KEY] = 値 \} \tabularnewline
構造化データを保存する & Roomの三種の神器：@Entity / @Dao / @Database \tabularnewline
ネットからデータを取る & Retrofitインターフェース＋suspend関数 \tabularnewline
通知を出す & NotificationChannel + NotificationCompat.Builder \tabularnewline
ログを出す & \texttt{Log.\allowbreak d("TAG", "message")} \tabularnewline
\hline
\end{tabularx}
\end{center}

\begin{center}\small\color{gray}Android開発クイックカード（2026年8月版）\end{center}

\booksection{よくある質問FAQ}

Q1：プログラミングがまったくできませんが、Android開発を学ぶにはどれくらいかかりますか？

本書の道順に従いましょう。第1〜2週で環境を整えてKotlinの基礎を学び、第3〜4週でComposeを学んで画面を作り、第5〜6週で家計簿プロジェクトを完成させ、第7〜8週で天気プロジェクトをやり遂げます。2か月で、ゼロから小さなアプリを独力で作れるまで進めます。条件は、毎日1〜2時間を投じて、本を読むだけでなく実際にコードを打つことです。

Q2：Javaも学ぶ必要がありますか？

不要です。新しいプロジェクト、新しいチュートリアル、新しい求人は、ほぼすべてKotlin基準です。大量のJavaに出会うのは「10年前のプロジェクトの保守」を頼まれたときだけです。Kotlinを終えてから、仕事で必要になれば、Javaの補習もすぐ済みます。両者の土台は同じ仮想マシンの世界だからです。

Q3：エミュレーター？　それとも実機？

開発やデバッグにはエミュレーターが便利です（スナップショット、複数機種）。でも検収は必ず実機で行いましょう。操作感、性能、本物のネットワーク、センサーは、エミュレーターでは代えられません。条件が許すなら、ミドルクラスの実機を1台テスト機として買うのが、いちばん費用対効果の高い投資です。

Q4：Android開発は仕事にありつきやすいですか？

2026年の相場はこうです。単に「画面が書ける」だけの初級ポストは競争が激しいものの、「Kotlinが分かる＋アーキテクチャができる＋リリース作品がある」開発者は、いまだに引っ張りだこです。学習中のあなたにとって、いちばんいい履歴書は、ストアにある自分で作ったアプリです。

Q5：エラーで先へ進めなくなったら、どうすればいいですか？

黄金の3ステップです。(1) エラーの最初の行を、そのままStackOverflowや検索エンジンで調べる。(2) 公式ドキュメントの該当ページを確認する。(3) それでもだめなら、「期待する動作＋実際の現象＋エラー全文＋関係するコード」をまとめて、AIやコミュニティに質問する。ほとんどの問題は、最初のステップで解決します。

\booksection{おすすめ学習リソース}

\begin{center}
\begin{tabularx}{\linewidth}{|X|X|X|}
\hline
\textbf{種類} & \textbf{リソース} & \textbf{おすすめの理由} \tabularnewline
\hline
公式 & developer.android.com & Androidの公式ドキュメント。ComposeやKotlinのチュートリアルもここにあります。第一の権威 \tabularnewline
公式 & Kotlin公式サイト (kotlinlang.org) & 言語の公式ドキュメント＋オンラインの演習場 \tabularnewline
公式 & Google Codelabs & 手取り足取りのステップ式実習コース。初心者へのやさしさは満点 \tabularnewline
コミュニティ & Stack Overflow（androidタグ） & エラーで助けを求める最初の場所 \tabularnewline
コミュニティ & 掘金 / CSDNのAndroidコーナー & 中国語の実践記事が多い。新しさの見極めに注意（まず日付を見る） \tabularnewline
動画 & Google Developers公式チャンネル & 毎年の開発者会議での新技術の解説 \tabularnewline
本書連動 & 第7部『AIプログラミング実戦編』（第37〜47章） & AIでAndroid開発を支援するための完全な心得 \tabularnewline
\hline
\end{tabularx}
\end{center}

\begin{center}\small\color{gray}Android開発の学習リソース一覧（2026年8月）\end{center}

\booksection{Gradleのよく使うコマンド速査}

ふだんはボタンで十分ですが、コマンドラインのほうが効率的なこともあります（特にAIツールやCIパイプラインと組み合わせるとき）。

\begin{center}
\begin{tabularx}{\linewidth}{|X|X|}
\hline
\textbf{コマンド} & \textbf{役割} \tabularnewline
\hline
\texttt{.\allowbreak /\allowbreak gradlew assembleDebug} & デバッグ版APKをビルドする（Windowsでは ./ を外す） \tabularnewline
\texttt{.\allowbreak /\allowbreak gradlew assembleRelease} & リリース版をビルドする（署名の設定が必要） \tabularnewline
\texttt{.\allowbreak /\allowbreak gradlew test} & ユニットテストを全部実行する \tabularnewline
\texttt{.\allowbreak /\allowbreak gradlew lint} & 静的チェック：潜在的なバグや行儀の悪い書き方を見つける \tabularnewline
\texttt{.\allowbreak /\allowbreak gradlew clean} & ビルド成果物を空にする（「再起動の大技」のコンパイル版） \tabularnewline
\texttt{.\allowbreak /\allowbreak gradlew :app: dependencies} & 依存ツリーを見る（バージョン競合の調査に便利な道具） \tabularnewline
\hline
\end{tabularx}
\end{center}

\begin{center}\small\color{gray}Gradleのよく使うコマンド速査\end{center}

\begin{tcbinfo}{ワンポイント：ビルドでおかしくなったら、まずclean}
「昨日は動いたのに今日は動かない」という不思議な出来事は、十中六七がビルドキャッシュの乱れです。./gradlew clean のあとにビルドし直せば、オカルト的な問題の大半は治ります。それでもだめならAndroid Studioを再起動しましょう——この技は File Invalidate Caches と併用するのがお約束です。

\end{tcbinfo}

\booksection{おわりに：最初のアプリからN個目のアプリへ}

\booksection{おわりに：最初のアプリからN個目のアプリへ}

振り返れば、あなたは「Activityが何かも知らない」状態から始まって、環境を整え、Kotlinを学び、Composeで画面を描き、アプリにデータベースとネットワークをつなぎ、さらにテスト、パッケージング、出品まで進みました。ゼロから一を生み出す、いちばん大変な一区間を、歩き終えたところです。

技術の世界は変化が速いものです。フレームワークは更新され、道具は世代交代し、AIはますます多くの「キーボードを叩く仕事」を引き受けていきます。でも、ひとつだけ色あせないものがあります。ぼんやりしたアイデアを明確なステップに分解し、それを一歩ずつ実現していく力です。道具は増幅器、考え方は地盤です。増幅器は取り替えられますが、地盤は動かせません。

あなたはもう「スマホが使える人」だけではありません。頭の中のアイデアを、数億台のAndroid端末で実際に動くアプリに変えられるのです。今日からキーボードで創造しよう——これがAndroid開発者のロマンです。

\booksection{参考文献}

\begin{enumerate}
  \item Android Developers. (2026). Android Developer Guides and Jetpack Compose Documentation. developer.android.com.
  \item Kotlin Foundation. (2026). Kotlin Language Documentation. kotlinlang.org.
  \item Android Open Source Project. (2026). Android Platform Version History and API Levels.
  \item Google Play. (2026). Play Console Help: Publishing and App Signing.
  \item AndroidX Release Notes. (2026). Room / DataStore / Navigation / WorkManager releases.
  \item Square. (2026). Retrofit Documentation. github.com/square/retrofit.
\end{enumerate}
